% disclaimers.tex — GLS-2026-001 rigour-without-infrastructure concept paper.
% Every id below is either (a) one of the ids emitted by
% `./cli/glosa claim disclaimers` against the three project claim cards
% (GLOSA-CC-20260904-0201/0202/0203, `disclaimers_emitted`, task B1) or (b) a
% template-mandatory id filled honestly for this paper. Ids are the semantic
% ids of the single disclaimer catalogue; no numeric D1-D12 scheme.

\begin{disclaimers}
  \item[D-STANDPOINT] See \Cref{sec:standpoint}. This paper is written from a
    social-enterprise-practitioner/citizen standpoint, not disciplinary
    authority; a reader who skips to this section should read every claim
    above as offered from that position and access level, never from
    institutional credential.
  \item[D-NONEXPERT] The author is not a substitute for a credentialed expert
    in citizen-science studies, philosophy of science, library/information
    science, or empirical AI-evaluation methodology --- the disciplines this
    paper draws on without claiming standing in. Medicine, religious/fiqh
    authority, law, engineering, anthropology, political science, and formal
    computer science beyond what is mechanically executed here are explicitly
    not claimed.
  \item[D-SCOPE] Not applicable in the empirical-population sense (this is a
    \textsc{CONCEPTUAL} genre paper, \Cref{sec:genre}). Where a source's own
    sample size or population matters to a comparison row
    (\Cref{tab:h1-lit,tab:h2-lit,tab:h3-lit}), it is stated inline in that
    row rather than in one blanket population statement.
  \item[D-NONCLAIM] This paper does \emph{not} claim: that H1's mechanism has
    been run as a study; that H2's conceptual-portability claim has been
    tested against real readers; that H3's seeded-error benchmark has been
    built or executed; that the 48 citation cards behind
    \Cref{tab:h1-lit,tab:h2-lit,tab:h3-lit} constitute a systematic or
    exhaustive review (\texttt{search\_mode: TARGETED\_SEARCH} /
    \texttt{SCOPING\_SEARCH}, never \texttt{SYSTEMATIC\_REVIEW}); that any
    finding here is a medical, legal, or religious determination; or that
    cross-vendor AI checking (I3) is equivalent to, or a substitute for, an
    I5 human review.
  \item[D-AIFILL] See \Cref{sec:ai-disclosure}. AI (the AI assistant, "the AI assistant" seat)
    drafted candidate hypotheses H1/H2/H3, the literature-review runs, the
    claim cards' \texttt{current\_evidence}, \texttt{retrieved\_tool\_evidence},
    \texttt{retained\_record\_route}, \texttt{model\_calibration\_assumption},
    \texttt{prompt\_system\_constraint}, and \texttt{decision\_policy} fields,
    and this paper's own prose and \LaTeX{} build. None of this is treated as
    evidence on its own say-so (\Cref{sec:ai-disclosure}).
  \item[D-TIER] The single highest tier claimed anywhere in this paper is
    \DrTier{} (declared bridge / human narrative synthesis over a specified
    mechanism), with one \FiniteDiag{} exception: the mechanically executed
    checks reported in \Cref{sec:demonstration} (schema validation, disclaimer
    emission, manifest gate reads) and this paper's own \LaTeX{} compile. It
    cannot go higher because none of H1/H2/H3 has been run as a study, and no
    claim card behind this paper has passed a class-5 (I5) independent human
    check.
  \item[D-INDEPENDENCE] Every citation card behind
    \Cref{tab:h1-lit,tab:h2-lit,tab:h3-lit} was checked at independence class
    $I3$ (decorrelated cross-vendor AI route, \texttt{claim\_match\_verified\_by:
    route:cross-vendor-ai}) for \texttt{claim\_match}, and mechanically ($I3$-adjacent,
    Crossref/OpenAlex lookup) for \texttt{metadata\_verified}. The three
    project claim cards themselves (\Cref{sec:claim-matrix}) sit at
    \texttt{independent\_check.status: NONE} --- drafted and self-checked by
    the authoring AI session, never yet reviewed by a route independent of
    that session (this is $I0$ for the cards' own content, not to be conflated
    with the $I3$ citation-level checks that feed them). No external or
    institutional validation is sought or treated as available as a lever
    that would make any claim here more true; an outside reader's input, if
    it comes, is optional-level evidence at whatever independence class it
    actually reaches, never a certification gate.
  \item[D-DVP-NOT-K2] A decorrelated-vendor AI review pass (DVP) is at most an
    $I3$ route. \textbf{DVP $\neq$ K2}: no combination of AI-only checks,
    however many vendors, raises this paper's claims to K2. This paper's
    K-state is \texttt{K0} (\Cref{sec:kstate}); no claim in it is asserted
    above K0/\DrTier{}.
  \item[D-LEGAL-NEQ-EPISTEMIC] None identified --- no license, permit, or
    regulatory approval is referenced by this paper's claims. The general
    rule still holds: were one ever cited, it would make an action lawful,
    never raise epistemic warrant or certify a claim as true.
  \item[D-NOT-DIAGNOSTIC] Not applicable --- this paper makes no health,
    clinical, or treatment recommendation of any kind.
  \item[D-TRANSLATION] Not applicable: per founder ruling
    (BBL-2026-09-04-099, \Cref{sec:blackbox}), this concept paper is
    published in English only, with no Thai edition. Every founder line
    quoted in \Cref{sec:blackbox} nonetheless stays Thai, verbatim,
    untranslated in place, with an English gloss alongside it --- the
    English-only rule governs the paper's own prose, not the Blackbox Note's
    raw lines. Within the underlying claim cards, the Thai statement fields
    themselves carry \texttt{translation\_status: reviewed}.
  \item[D-DISSENT-PRESERVED] None recorded to date for H1/H2/H3 specifically.
    The literature-review dialogue tables (\Cref{tab:h1-lit,tab:h2-lit})
    already carry the closest real dissent this project has found against its
    own hypotheses (the Mayes 2022 boundary-work finding against H1
    \citep{cite-cstp-boundary-work-001}; the Scholarly Kitchen credentialism
    case against H2 \citep{cite-rigour-without-infrastructure-004}); these are
    stated as open challenges in \Cref{sec:literature}, not resolved or
    smoothed over.
  \item[D-COMPARISON] This paper does not claim priority or precedence.
    \Cref{sec:literature} states every source relation as \emph{agrees},
    \emph{disagrees}, or \emph{orthogonal/undetermined} --- never as a
    priority contest, never as ``no other work does X.''
  \item[D-CITATION-UNVERIFIED] None identified for the 48 citation cards cited
    in \Cref{tab:h1-lit,tab:h2-lit,tab:h3-lit} --- every one carries
    \texttt{status: VERIFIED} (\texttt{metadata\_verified: true},
    \texttt{claim\_match\_verified: true}) in its litreview manifest as of
    2026-09-04; none is cited from memory (\Cref{sec:demonstration},
    BBL-2026-09-04-100/101). Where a source's own full text was not fetched
    (e.g.\ \texttt{fetch\_status: PAYWALLED\_ABSTRACT\_ONLY}), the citation
    card's \texttt{scope} field names the resulting evidentiary weakness
    (\texttt{PARAPHRASE} rather than \texttt{DIRECT\_QUOTATION}) inline in
    the relevant table row rather than silently upgrading it.
  \item[D-BLACKBOX-NOTE] This paper carries the mandatory Blackbox Note
    appendix (\Cref{sec:blackbox}). Note id: \texttt{BB-2026-09-04-01}
    (\texttt{projects/GLS-2026-001\_rigour-without-infrastructure/blackbox\_note.yaml});
    curated public source:
    \texttt{blackbox/BLACKBOX\_NOTE\_glosa-paper\_2026-09-04.md}; public log
    entries \texttt{BBL-2026-09-04-083, -086, -088, -089, -091, -095, -097,
    -098, -099, -100, -101} (\texttt{blackbox/log/entries.jsonl}, concept DOI
    \texttt{10.5281/zenodo.22302518}).
  \item[D-K-STATE] Overall K-state: \texttt{K0} (internal working note). Every
    load-bearing claim in this paper --- H1, H2, H3, and every finding in
    \Cref{sec:demonstration} --- shares this same K0 state; there is no
    per-claim exception. This matches the three underlying project claim
    cards' own \texttt{k\_state: K0} (\Cref{sec:claim-matrix}).
  \item[D-AUTHORSHIP] See \Cref{sec:ai-disclosure}. The problem, the research
    question, and the selection among H1/H2/H3 are the founder's own act
    (\Cref{sec:standpoint-and-ownership}); AI drafted every candidate,
    citation search, and this paper's prose, disclosed per claim
    (\texttt{produced\_by}, \Cref{tab:claim-matrix}), and never substitutes
    for the author's judgment or responsibility
    (\textbf{AIContribution $\neq$ EpistemicResponsibility}).
  \item[D-REVISION-LIVE] The three underlying claim cards
    (\texttt{GLOSA-CC-20260904-0201/0202/0203}) are at revision 2 of an
    actively edited project; \texttt{status: Draft} and
    \texttt{independent\_check.status: NONE} on all three reflect that the
    founder has not yet signed off on any specific card instance, even though
    the founder has already ruled on genre (BBL-097) and hypothesis selection
    (BBL-095) at the project level.
  \item[D-NO-EPISTEMIC-VETO] No agency --- including any AI vendor, reviewer,
    or institution named in this paper --- holds an epistemic veto over H1,
    H2, or H3 merely by title, credential, or infrastructure ownership; any
    agency may propose, use, test, challenge, fork, translate, revise,
    restrict, or withdraw this work without needing permission, provided
    lineage is preserved (per the repo's own knowledge-validation stance,
    \texttt{CLAUDE.md}/\texttt{AGENTS.md}).
  \item[D-CANDIDATE-STATUS] H1, H2, and H3 are candidate hypotheses drafted by
    AI; the act of selecting among them (all three, per BBL-2026-09-04-095)
    was the founder's alone (\Cref{sec:standpoint-and-ownership}). No claim
    card here asserts candidate status was itself a form of testing.
  \item[D-EXTERNAL-INPUT] The lens this paper's readout-translation and
    hypothesis-drafting steps were performed under is named, cited, and
    linked throughout (\Cref{sec:lens}); naming it credits its source and
    lets a reader replace it --- it does not make the lens, or any hypothesis
    drafted under it, true.
\end{disclaimers}
